% ==============================================================================
% PLANTILLA MAESTRA GENERAL - UCNL
% ==============================================================================
% Uso:
% 1. Duplica este archivo y renombra la copia por materia, semana y actividad.
% 2. Actualiza documenttitle, documentsubtitle, documentsubject, coursename,
%    coursecode, universitydepartment y datos de la figura docente.
% 3. Lee la planeacion semanal y NO pegues las instrucciones completas dentro del
%    producto final. Usa la planeacion como lista de cumplimiento.
% 4. Revisa la bibliografia indicada, toma notas, registra citas y redacta con
%    comprension propia.
% 5. Identifica el producto academico solicitado y construyelo dentro del
%    documento siempre que sea posible.
% 6. Entrega en PDF con la nomenclatura indicada en el aula.
%
% Formato institucional recomendado:
% - Fuente tipo Helvetica, equivalente visual a Arial.
% - Interlineado 1.5.
% - Citas autor-anio con natbib: usar \citep{clave} o \citet{clave}.
% - Referencias en bibliografia BibTeX, formato APA 7 en cuerpo y referencias.
% - Caratula, resumen breve, desarrollo, producto solicitado, conclusion y
%   bibliografia.
% - Redaccion academica clara, coherente, cohesionada, formal y sin faltas.
% - La portada puede llevar una marca de agua institucional discreta; ajustar
%   las variables coverwatermark... si se requiere apagarla o cambiar intensidad.
%
% Criterios generales de cumplimiento:
% - Responder exactamente a la actividad solicitada en la planeacion.
% - Identificar conceptos, elementos, criterios o categorias de forma precisa.
% - Analizar, interpretar y tomar postura; no solo copiar fuentes.
% - Organizar la informacion por jerarquia, secuencia y criterios claros.
% - Incluir conclusion personal cuando la actividad lo permita o lo solicite.
% - Usar citas y referencias en APA 7.
% - Evitar copiar las indicaciones de la actividad en el cuerpo final.
%
% Regla para productos visuales:
% - Si la actividad pide tabla, cuadro comparativo, esquema, mapa conceptual,
%   mapa mental, linea de tiempo, diagrama, matriz, lista de verificacion, red
%   conceptual, SQA, organizador grafico o producto similar, realizarlo
%   preferentemente con LaTeX/TikZ para mantener nitidez, evidencia editable y
%   diseno controlado.
% - Si la consigna exige Canva, Miro u otra herramienta externa, insertar una
%   captura nitida y conservar la explicacion dentro del documento.
% - Para tablas extensas, usar pagina limpia, sin encabezados ni pies, landscape,
%   letra pequena y restaurar el estilo al terminar.
%
% Criterios de redaccion personal:
% - No escribir para llenar espacio; escribir para sostener una idea.
% - Convertir informacion en criterio: que significa, como funciona, que problema
%   resuelve, que limite presenta y que consecuencia produce.
% - Estructura mental recomendada: problema -> sistema -> consecuencia -> postura.
% - Usar pensamiento de diagnostico: entrada -> proceso -> salida.
% - Cada parrafo debe tener funcion: presentar, explicar, comparar, valorar o
%   concluir.
% - Usar primera persona con moderacion academica: "Considero que",
%   "Desde esta perspectiva", "La posicion que sostengo es". Evitar "yo creo".
%
% Notas de enfoque personal segun enfoque_personal.txt:
% - La voz debe ser academica, tecnica, directa, estrategica y funcional.
% - El texto no debe sonar como estudiante pasivo que recopila: debe sonar como
%   alguien que interpreta, ordena, cruza datos y busca consecuencia practica.
% - Antes de redactar, identificar la funcion de cada seccion: abrir problema,
%   definir criterio, demostrar con evidencia, comparar, valorar o concluir.
% - Evitar frases genericas como "es importante porque si" o "desde siempre".
%   Sustituirlas por relaciones verificables: causa, efecto, limite, costo,
%   funcion, riesgo, implicacion o mejora.
% - Cada concepto debe pasar de definicion a analisis. Formula util:
%   concepto -> funcion -> riesgo si se aplica mal -> consecuencia -> postura.
% - La postura personal no debe ser opinion suelta. Debe tener tres piezas:
%   posicion, razon y efecto. Ejemplo: "Considero que X, porque Y; por ello Z".
% - Integrar experiencia o criterio propio sin perder formalidad: no narrar la
%   vida personal, sino usar mirada practica para explicar por que el tema
%   importa en educacion, derecho, instituciones, comunicacion o profesion.
%
% Estilo observado en actividades ya realizadas:
% - Abrir con una tesis concreta:
%   "El Derecho no solo ordena conductas: organiza poder, define limites y
%   distribuye consecuencias."
% - Formular el problema en negativo y positivo:
%   "El problema no es solo..., sino..."
% - Convertir el producto visual en argumento:
%   "El cuadro no solo recopila datos; reconstruye la funcion, el problema y el
%   punto que requiere correccion."
% - Cerrar con aprendizaje transferible:
%   "No basta con saber que dice la norma; tambien debe analizarse a quien
%   protege, a quien excluye y que razones permiten defenderla publicamente."
%
% Nivel cognitivo esperado:
% - Recordar: usar solo para definir conceptos basicos. No debe ser el nivel
%   dominante en una actividad ponderable.
% - Comprender: explicar con palabras propias y relacionar con el tema.
% - Aplicar: llevar el concepto a un caso, norma, texto, correo, dilema o hecho.
% - Analizar: separar elementos, detectar causas, efectos, tensiones y limites.
% - Evaluar: sostener una postura con criterios, fuentes y consecuencias.
% - Crear: elaborar producto propio: mapa, cuadro, linea de tiempo, ensayo,
%   propuesta, reformulacion, glosario, portafolio o solucion argumentada.
%
% Regla de nivel minimo:
% - Si la actividad vale puntos y pide producto, no quedarse en recordar o
%   comprender. Debe llegar por lo menos a aplicar y analizar; si pide conclusion
%   o postura, debe llegar a evaluar.
%
% Uso de las 100 tecnicas didacticas:
% - Esta plantilla conserva el manual "Tecnicas didacticas del aprendizaje" que
%   se incorporo en las plantillas de Filosofia del Derecho, Etica y Moral
%   juridica, y Redaccion en contextos virtuales.
% - Cuando la planeacion mencione una tecnica de la coleccion UnADM, identificar:
%   definicion de la tecnica, producto esperado, pasos de construccion, criterios
%   de revision y cierre reflexivo.
% - La tecnica didactica no debe verse como adorno. Debe mostrar seleccion,
%   jerarquia, relacion entre ideas y lectura propia.
% - Siempre que sea viable, construir la tecnica en LaTeX/TikZ: cuadro
%   comparativo, mapa conceptual, esquema, mapa mental, linea de tiempo, matriz,
%   red conceptual, diagrama de flujo, tabla SQA o lista de verificacion.
%
% Checklist antes de entregar:
% [ ] El titulo, materia, semana, docente y fecha estan actualizados.
% [ ] El objetivo coincide con la planeacion.
% [ ] El producto solicitado aparece visible, rotulado y completo.
% [ ] La introduccion plantea un problema o postura, no solo anuncia el tema.
% [ ] El desarrollo explica, conecta y analiza.
% [ ] Hay citas APA autor-anio dentro del texto.
% [ ] La bibliografia final coincide con las fuentes citadas.
% [ ] La conclusion responde que se comprendio y que postura se sostiene.
% [ ] El nivel cognitivo llega al menos a aplicacion y analisis.
% [ ] Si hay conclusion/postura, incluye posicion, razon y consecuencia.
% [ ] La tecnica didactica solicitada esta construida y explicada.
% [ ] No aparecen instrucciones copiadas de la planeacion dentro del producto.
% [ ] Ortografia, acentos, sintaxis y nombres propios revisados.
% [ ] Si se uso IA como apoyo, el texto fue comprendido, revisado y apropiado.
% ==============================================================================

% CREACION DEL DOCUMENTO
\documentclass[
	spanish,
	letterpaper, oneside
]{article}

% ==============================================================================
% INFORMACION DEL DOCUMENTO
% Actualizar en cada actividad.
% Ejemplos:
% \def\documenttitle {Cuadro comparativo de conceptos fundamentales}
% \def\documentsubtitle {Actividad 1 - Administracion de operaciones}
% \def\documentsubject {Actividad 1}
% ==============================================================================
\def\documenttitle {Actividad 1 - Administracion de operaciones}
\def\documentsubtitle {Actividad 1 - Administracion de operaciones}
\def\documentsubject {Actividad 1}

\def\documentauthor {Martin Jonathan de la Cruz}
\def\coursename {Administracion de operaciones}
\def\coursecode {ADM-OPE-T8}

\def\universityname {Universidad Ciudadana de Nuevo Leon}
\def\universityfaculty {Programa academico UCNL}
\def\universitydepartment {Administracion de operaciones}
\def\universitydepartmentimage {departamentos/logo-ucnl}
\def\universitydepartmentimagecfg {height=1.57cm}
\def\universitylocation {Monterrey, Nuevo Leon}

% INTEGRANTES, PROFESORES Y FECHAS
\def\authortable {
	\begin{tabular}{ll}
		\textbf{Alumno:}
		& \begin{tabular}[t]{l}
			 Martin Jonathan de la Cruz \\
		\end{tabular} \\
		Matricula: & UCNL \\

		\textit{Figura docente:}
		& \begin{tabular}[t]{l}
			Nombre \\ Apellido
		\end{tabular} \\
		& \\
		\multicolumn{2}{l}{Fecha de realizacion: \today} \\
		\multicolumn{2}{l}{\universitylocation}
	\end{tabular}
}

% IMPORTACION DEL TEMPLATE BASE
\input{template}

% FORMATO GENERAL
\usepackage{helvet}
\renewcommand{\familydefault}{\sfdefault}

% TABLAS Y PRODUCTOS VISUALES
\usepackage{array}
\usepackage{longtable}
\usepackage{pdflscape}
\usepackage{tikz}
\usetikzlibrary{
	arrows.meta,
	positioning,
	calc,
	matrix,
	fit,
	shapes.geometric,
	shadows.blur
}

% CITAS EN FORMATO AUTOR-ANIO, COMPATIBLE CON APA 7 EN EL CUERPO DEL TEXTO
% Usar:
% - Cita parentetica: \citep{clave}
% - Cita narrativa: \citet{clave}
% - Varias fuentes: \citep{clave1,clave2}
\setcitestyle{authoryear,open={(},close={)}}

% ==============================================================================
% AYUDAS DE REDACCION
% Quitar o comentar \pendiente cuando el documento final este terminado.
% ==============================================================================
\newcommand{\pendiente}[1]{\textcolor{red}{Refuerzo Ciclo A materializado}}

% ==============================================================================
% MARCA DE AGUA EN PORTADA
% - Se aplica solo a la portada porque usa \AddToShipoutPictureBG*.
% - Para apagarla en alguna actividad: \def\coverwatermarkenabled {false}
% - Usar opacidad baja para que no compita con titulo, datos ni logotipo.
% ==============================================================================
\def\coverwatermarkenabled {true}
\def\coverwatermarkimage {img/departamentos/logo-nuevo-leon.png}
\def\coverwatermarkopacity {0.16}
\def\coverwatermarkwidth {0.72\paperwidth}
\def\coverwatermarkxshift {0cm}
\def\coverwatermarkyshift {-0.35cm}

\newcommand{\insertcoverwatermark}{%
	\ifthenelse{\equal{\coverwatermarkenabled}{true}}{%
		\AddToShipoutPictureBG*{%
			\begin{tikzpicture}[remember picture,overlay]
				\node[
					opacity=\coverwatermarkopacity,
					inner sep=0pt,
					xshift=\coverwatermarkxshift,
					yshift=\coverwatermarkyshift
				] at (current page.center) {%
					\includegraphics[width=\coverwatermarkwidth]{\coverwatermarkimage}%
				};
			\end{tikzpicture}%
		}%
	}{}%
}

\begin{document}

% PORTADA
\insertcoverwatermark
\templatePortrait

% CONFIGURACION DE PAGINA Y ENCABEZADOS
\templatePagecfg
\onehalfspacing

% ==============================================================================
% RESUMEN / ABSTRACT
% Debe ser breve. Formula recomendada:
% Tema + objetivo + tecnica/producto + enfoque personal.
% ==============================================================================
\begin{abstractd}
	\textbf{Objetivo:} \textit{Refuerzo Ciclo A: Redactar el objetivo con verbo academico: analizar, identificar, aplicar, comparar, explicar, evaluar, interpretar, argumentar o proponer. Debe coincidir con la planeacion semanal.}

	\textbf{Producto elaborado:} \textit{Refuerzo Ciclo A: Indicar el producto solicitado: cuadro comparativo, mapa conceptual, linea de tiempo, estudio de caso, lista de verificacion, reporte, ensayo, foro, portafolio, glosario, resumen, esquema u otro.}

	\textbf{Enfoque de analisis:} \textit{Refuerzo Ciclo A: Enunciar la postura central o criterio de lectura. Debe mostrar desde donde se analizara el tema y por que importa.}

	\textbf{Nivel cognitivo:} \textit{Refuerzo Ciclo A: Indicar el nivel dominante de la actividad: comprender, aplicar, analizar, evaluar o crear. Evitar que el trabajo se quede solo en recordar informacion.}
\end{abstractd}

% TABLA DE CONTENIDOS - INDICE
\templateIndex

% CONFIGURACIONES FINALES
\templateFinalcfg

% ======================= INICIO DEL DOCUMENTO =======================

% ==============================================================================
% SECCION 1: INTRODUCCION
% Apertura recomendada:
% - Iniciar con una afirmacion clara, no con una frase generica.
% - Traducir la consigna a un problema academico, juridico, etico, comunicativo
%   o profesional, segun la materia.
% - Explicar por que importa.
% - Cerrar anticipando el enfoque personal.
% ==============================================================================
\section{Introduccion}

La administración de operaciones no se limita a ejecutar tareas internas; su función real es transformar recursos en resultados con criterios de costo, tiempo, calidad y capacidad de respuesta. El problema central que aborda esta actividad es que, en muchos entornos, las decisiones operativas se toman de forma aislada y reactiva, lo que provoca ineficiencias acumuladas: retrabajos, demoras, sobrecostos y pérdida de valor para el usuario final. En ese sentido, el reto no es únicamente describir conceptos, sino comprender cómo se articulan para sostener desempeño operativo.

El objetivo de este reporte es analizar los componentes clave de la operación desde una lógica de diagnóstico (entrada, proceso y salida), para identificar cómo cada decisión técnica impacta el cumplimiento de metas organizacionales. El enfoque de trabajo parte de una premisa: una operación bien diseñada no es la que produce más, sino la que produce de manera estable, medible y alineada con la demanda. Bajo este criterio, el desarrollo se orienta a interpretar relaciones causales y no solo a enumerar definiciones, de modo que el producto final refleje comprensión aplicada y capacidad de evaluación.

\section{Desarrollo}

\subsection{Contexto y problema}

\textit{Refuerzo Ciclo A: Presentar el contexto de la actividad y transformar la consigna en problema de analisis.}

\subsection{Marco conceptual}

\textit{Refuerzo Ciclo A: Explicar los conceptos centrales con fuentes. No acumular definiciones: cada concepto debe mostrar funcion, limite o consecuencia.}

\subsection{Nivel cognitivo aplicado}

\textit{Refuerzo Ciclo A: Explicar brevemente que operacion cognitiva exige la actividad: comparar, diagnosticar, clasificar, interpretar, argumentar, evaluar o crear. Relacionar esa operacion con el producto solicitado.}

\subsection{Producto solicitado}

\textit{Refuerzo Ciclo A: Insertar aqui el producto de la actividad. Si es visual, realizarlo en LaTeX/TikZ siempre que sea posible.}

\subsection{Tecnica didactica aplicada}

\pendiente{Cuando la planeacion indique una tecnica de las 100 Tecnicas Didacticas de Ensenanza y Aprendizaje, explicar aqui como se uso. Ejemplo: "La tecnica seleccionada organiza la informacion mediante criterios visibles, jerarquia y cierre reflexivo; por ello, el producto no solo presenta datos, sino que demuestra analisis y comprension" \citep{unadm100TecnicasDidacticas2023}.}

% ==============================================================================
% BLOQUE OPCIONAL: TABLA O CUADRO COMPARATIVO AMPLIO
% Usar cuando el producto necesite mucho ancho:
% - \clearpage para iniciar en pagina limpia.
% - \pagestyle{empty} para quitar encabezado y pie en el producto amplio.
% - landscape para aprovechar el ancho horizontal.
% - scriptsize, tabcolsep y arraystretch para controlar legibilidad.
% - Restaurar \pagestyle{fancy} al terminar.
%
% Si la tabla es corta, puede omitirse landscape y usarse longtable normal.
% ==============================================================================
% \clearpage
% \pagestyle{empty}
% \begin{landscape}
% \begingroup
% \scriptsize
% \setlength{\tabcolsep}{4pt}
% \renewcommand{\arraystretch}{1.35}
% \begin{longtable}{@{}p{0.18\linewidth}p{0.39\linewidth}p{0.39\linewidth}@{}}
% \caption{Titulo del cuadro comparativo}\label{tab:cuadro-comparativo}\\
% \toprule
% \textbf{Criterio} & \textbf{Elemento 1} & \textbf{Elemento 2} \\
% \midrule
% \endfirsthead
% \toprule
% \textbf{Criterio} & \textbf{Elemento 1} & \textbf{Elemento 2} \\
% \midrule
% \endhead
% Concepto & \textit{Refuerzo Ciclo A: Contenido sintetico con cita si aplica.} & \textit{Refuerzo Ciclo A: Contenido sintetico con cita si aplica.} \\
% Caracteristicas & \textit{Refuerzo Ciclo A: Contenido.} & \textit{Refuerzo Ciclo A: Contenido.} \\
% Funcion & \textit{Refuerzo Ciclo A: Contenido.} & \textit{Refuerzo Ciclo A: Contenido.} \\
% Valoracion critica & \textit{Refuerzo Ciclo A: Contenido.} & \textit{Refuerzo Ciclo A: Contenido.} \\
% \bottomrule
% \end{longtable}
% \endgroup
% \end{landscape}
% \pagestyle{fancy}

% ==============================================================================
% BLOQUE OPCIONAL: PRODUCTO VISUAL EN TIKZ
% Usar para mapa conceptual, mapa mental, esquema, flujo, proceso o diagrama.
% Cada flecha debe expresar relacion; no usar nodos sueltos sin conectores.
% ==============================================================================
% \begin{figure}[H]
% 	\centering
% 	\begin{tikzpicture}[
% 		node distance=1.3cm,
% 		box/.style={rectangle, rounded corners=2pt, draw=black!65, fill=black!4, align=center, minimum width=3.2cm, minimum height=0.9cm},
% 		arrow/.style={-{Latex[length=2.5mm]}, thick, draw=black!65}
% 	]
% 		\node[box] (centro) {Concepto central};
% 		\node[box, right=of centro] (relacion) {Concepto relacionado};
% 		\node[box, below=of relacion] (efecto) {Consecuencia};
%
% 		\draw[arrow] (centro) -- node[above]{implica} (relacion);
% 		\draw[arrow] (relacion) -- node[right]{produce} (efecto);
% 	\end{tikzpicture}
% 	\caption{Titulo del producto visual.}
% \end{figure}

\subsection{Analisis propio}

El producto elaborado evidencia que el desempeño operativo depende menos de acciones aisladas y más de la coherencia del sistema completo. El hallazgo principal es que variables como capacidad, estandarización, secuencia del proceso y control de calidad funcionan como un conjunto interdependiente: cuando una de ellas se desajusta, el efecto no queda local, sino que se propaga al flujo total. Esta lectura permite pasar de la descripción a la interpretación técnica.

Desde este marco, el criterio que sostengo es funcional: toda decisión en operaciones debe evaluarse por su efecto en flujo, variabilidad y cumplimiento. Su función es reducir incertidumbre operativa y hacer predecible el resultado. Si ese criterio no se aplica, la consecuencia directa es operar con aparente actividad pero con baja productividad real; se trabaja más, pero se entrega peor o más tarde. En términos administrativos, esto deteriora indicadores clave y limita la capacidad de mejora continua.

La consecuencia estratégica de lo anterior es clara: sin una lógica de integración entre planeación, ejecución y control, la operación pierde ventaja competitiva aunque existan recursos suficientes. Por ello, mi postura técnica es que la administración de operaciones debe tratarse como sistema de decisiones basado en evidencia, donde cada ajuste se justifique por su impacto en desempeño global y no por conveniencia inmediata del área.

\section{Postura personal}

Considero que la administración de operaciones es una competencia directiva central, porque conecta la estrategia con resultados verificables en tiempo, costo y calidad. La razón es que, en la práctica, los objetivos institucionales solo se materializan cuando los procesos están diseñados, medidos y corregidos con criterios consistentes. La consecuencia de asumir esta perspectiva es que la gestión deja de depender de respuestas improvisadas y se orienta a decisiones trazables, con mayor control de riesgos y mejor aprovechamiento de recursos.

Desde mi posición, el aprendizaje relevante de esta actividad es que operar bien no equivale a mantener ocupados los recursos, sino a generar valor sostenido con estabilidad y mejora progresiva. Por ello, sostengo una postura de gestión operativa preventiva: diagnosticar antes de corregir, estandarizar antes de escalar y evaluar impacto antes de implementar cambios.

% --- Ciclo A: refuerzo editorial materializado ---
\section*{Refuerzo editorial Ciclo A}
Este apartado consolida la mejora editorial incremental del nodo a partir del estándar maduro de AulaTeX. El refuerzo atiende brechas detectadas de bibliografía, citas, análisis propio, transferencia y cierre argumentado, sin sustituir la consigna original ni copiar contenido de los casos maestros.

\subsection*{Tesis de trabajo}
La actividad debe sostener una postura verificable: el producto solicitado no sólo organiza información, sino que permite interpretar el problema disciplinar, justificar decisiones y reconocer sus consecuencias académicas o profesionales.

\subsection*{Análisis propio}
El hallazgo central es que una entrega madura requiere pasar de la descripción a la interpretación. Por ello, el desarrollo debe explicar qué revela el producto construido, por qué es relevante para la materia y qué criterio permite evaluar su pertinencia. Esta operación convierte la evidencia reunida en argumento académico.

\subsection*{Transferencia profesional}
La transferencia consiste en vincular el producto con situaciones reales de desempeño: diagnóstico, toma de decisiones, comunicación de resultados, cumplimiento normativo, intervención educativa o resolución de problemas según el campo disciplinar. Así, la actividad adquiere valor formativo más allá de la plantilla.

\subsection*{Cierre argumentado}
En consecuencia, la calidad de la actividad depende de que el producto visible, las fuentes y la conclusión formen una secuencia coherente. La conclusión debe recuperar la tesis, indicar el principal aprendizaje y señalar una implicación práctica o conceptual.

% Brechas locales atendidas por Ciclo A: bibliografia_insuficiente, pendientes_editoriales
% Reglas globales aplicadas:
% - Priorizar cierre de brecha recurrente: bibliografia_insuficiente (427 nodos).
% - Priorizar cierre de brecha recurrente: conclusion_argumentada (375 nodos).
% - Priorizar cierre de brecha recurrente: citas_insuficientes (328 nodos).
% - Priorizar cierre de brecha recurrente: analisis_propio (288 nodos).
% --- Fin Ciclo A ---

\section{Conclusion}

En conclusión, esta actividad permitió comprender que la administración de operaciones es un sistema de coordinación y control que traduce recursos en resultados medibles, y no solo un conjunto de tareas de ejecución. Se identificó que el núcleo del problema operativo está en la falta de integración entre decisiones de capacidad, proceso y calidad, lo cual produce ineficiencias acumuladas y afecta el cumplimiento organizacional.

Con base en el análisis realizado, la postura final que sostengo es que la efectividad operativa depende de decisiones técnicas articuladas, evaluadas por su consecuencia en el desempeño global del sistema. En ese sentido, el aprendizaje obtenido no se limita al plano conceptual: aporta un criterio de actuación profesional para diagnosticar, priorizar y mejorar procesos con enfoque en valor, estabilidad y mejora continua.

\bibliography{administracion-de-operaciones}

% FIN DEL DOCUMENTO

\subsection*{Citas de refuerzo Ciclo A}
La mejora editorial se vincula con la memoria documental disponible mediante las referencias \citep{unadm100TecnicasDidacticas2023,cicloARefuerzo02,cicloARefuerzo03,cicloARefuerzo04,cicloARefuerzo05}. Estas citas deben revisarse en una etapa disciplinar fina para sustituir o complementar fuentes genéricas por literatura específica del tema.

\end{document}





